% Part: counterfactuals
% Chapter: introduction
% Section: strict-conditional

\documentclass[../../../include/open-logic-section]{subfiles}

\begin{document}

\olfileid{cnt}{int}{str}

\olsection{કડક શરતી વિધાન}

લૂઈસે \emph{કડક શરતી વિધાન}~$\strictif$ રજૂ કર્યું અને દલીલ કરી કે
ભૌતિક શરતી વિધાનને બદલે આ વિધાન ફલિતતાને અનુરૂપ છે. સત્યલક્ષી મોડલ
તર્કશાસ્ત્રમાં $!A \strictif !B$ને $\Box(!A \lif
!B)$ તરીકે વ્યાખ્યાયિત કરી શકાય છે. એટલે અનુરૂપ ભૌતિક શરતી
વિધાન અનિવાર્ય હોય ત્યારે અને માત્ર ત્યારે કડક શરતી વિધાન (કોઈ
જગતમાં) સત્ય હોય છે.

ભૌતિક શરતી વિધાનના વિરોધાભાસોની સામે કડક શરતી વિધાન કેટલું સફળ થાય
છે? અસત્ય પૂર્વપક્ષ ધરાવતું કડક શરતી વિધાન હોય કે સત્ય પરિણામ
ધરાવતું હોય, તે સત્ય પણ હોઈ શકે અને અસત્ય પણ. વધુમાં,
$(!A \strictif !B) \lor (!B \strictif !A)$ માન્ય નથી. કડક શરતી
વિધાન $!A \strictif !B$, $\lnot !A \lor !B$ને સમતુલ્ય પણ નથી;
એટલે તે સત્યતાફલનલક્ષી નથી.

અહીં \Log{S5}ના પરિપ્રેક્ષ્યમાં નીચેના સંબંધો મળે છે:
\sourcecorrection{OLMD-107}{આ ફલિતતા-સંબંધો તથા નીચેની અનિવાર્યતા-દલીલો દરેક મોડલ તંત્રમાં માન્ય નથી; અહીં S5નો પરિપ્રેક્ષ્ય સ્પષ્ટ કર્યો છે.}
\begin{align}
  !A \strictif !B & \Entails \lnot !A \lor !B
  \text{ પરંતુ:}\\
  \lnot !A \lor !B & \Entails/ !A \strictif !B\\
  !B & \Entails/ !A \strictif !B\\
  \lnot !A & \Entails/ !A \strictif !B\\
  \lnot(!A \strictif !B) & \Entails/ !A \land \lnot !B
  \text{ પરંતુ:}\\
  !A \land \lnot !B & \Entails \lnot(!A \strictif !B)
  \intertext{તેમ છતાં કડક શરતી વિધાન પર મોડસ પોનેન્સનો નિયમ લાગુ પડે છે:}
  !A, !A \strictif !B & \Entails !B
  \intertext{કડક શરતી વિધાન પરિણામોને એકત્ર કરે છે:}
  !A \strictif !B,  !A \strictif !C & \Entails !A \strictif (!B \land !C)
  \intertext{કડક શરતી વિધાનમાં પૂર્વપક્ષને મજબૂત કરી શકાય છે:}
  !A \strictif !B & \Entails (!A \land !C) \strictif !B
  \intertext{કડક શરતી વિધાન સંક્રામક પણ છે:}
  !A \strictif !B, !B \strictif !C & \Entails !A \strictif !C
  \intertext{છેવટે કડક શરતી વિધાન તેના પ્રતિપક્ષી સ્વરૂપ સાથે સમતુલ્ય છે:}
  !A \strictif !B & \Entails \lnot !B \strictif \lnot !A\\
  \lnot !B \strictif \lnot !A & \Entails !A \strictif !B
\end{align}
\sourcecorrection{OLMD-108}{મૂળમાં પાંચમી ફલિતતા-પંક્તિએ ભૌતિક સંયોજક મૂક્યો છે; તે રૂપ માટે દર્શાવેલી અ-ફલિતતા ખોટી છે, તેથી અહીં કડક સંયોજક મૂક્યો છે.}

\begin{prob}
  કડક શરતી વિધાન માટે માન્ય ન હોય એવા ફલિતતા-સંબંધોના
  \Log{S5}-પ્રત્યુદાહરણો આપો, એટલે કે નીચેનાં માટે:
  \begin{enumerate}
  \item $\lnot p \Entails/ \Box(p \lif q)$
  \item $q \Entails/ \Box(p \lif q)$
  \item $\lnot\Box(p \lif q) \Entails/ p \land \lnot q$
  \item $\Entails/ \Box(p \lif q) \lor \Box(q \lif p)$
  \end{enumerate}
\end{prob}

\begin{prob}
  માન્ય ફલિતતા-સંબંધો કડક શરતી વિધાન માટે સાચા છે તે નીચેના
  \Log{S5}-પ્રમાણો આપીને બતાવો:
  \begin{enumerate}
  \item  $\Box(!A \lif !B) \Entails \lnot !A \lor !B$
  \item $!A \land \lnot !B \Entails \lnot\Box(!A \lif !B)$
  \item $!A, \Box(!A \lif !B) \Entails !B$
  \item $\Box(!A \lif !B), \Box(!A \lif !C) \Entails \Box(!A \lif (!B
    \land !C))$
  \item $\Box(!A \lif !B) \Entails \Box((!A \land !C) \lif !B)$
  \item $\Box(!A \lif !B), \Box(!B \lif !C) \Entails \Box(!A
    \lif !C)$
  \item $\Box(!A \lif !B) \Entails \Box(\lnot !B \lif \lnot !A)$
  \item $\Box(\lnot !B \lif \lnot !A) \Entails \Box(!A \lif !B)$
  \end{enumerate}
  \end{prob}

તેમ છતાં કડક શરતી વિધાનના પોતાના ``વિરોધાભાસો'' છે. અસત્ય પૂર્વપક્ષ
અથવા સત્ય પરિણામ ધરાવતું ભૌતિક શરતી વિધાન સત્ય હોય છે; તેવી જ રીતે
\emph{અનિવાર્યપણે} અસત્ય પૂર્વપક્ષ અથવા અનિવાર્યપણે સત્ય પરિણામ ધરાવતું
કડક શરતી વિધાન સત્ય હોય છે. વધુમાં, \Log{S5}માં દરેક સત્ય કડક શરતી
વિધાન અનિવાર્યપણે સત્ય છે અને દરેક અસત્ય કડક શરતી વિધાન અનિવાર્યપણે
અસત્ય છે. બીજા શબ્દોમાં,
\begin{align}
  \Box \lnot !A & \Entails !A \strictif !B\\
  \Box !B & \Entails !A \strictif !B\\
  !A \strictif !B& \Entails \Box(!A \strictif !B)\\
  \lnot(!A \strictif !B) & \Entails \Box\lnot(!A \strictif !B)
\end{align}
$\strictif$ને ``ફલિત કરે છે'' એમ સમજીએ તો આ કોઈ સમસ્યા નથી.
આખરે તાર્કિક ફલિતતા-સંબંધો ગાણિતિક તથ્યો છે; તેથી તે સંજોગાધીન
હોઈ શકતા નથી. પરંતુ $\strictif$ને ``જો \dots તો \dots''નો અર્થ
પકડતો તાર્કિક સંયોજક માનવો હોય, તો ખાસ કરીને છેલ્લા બે સંબંધો
સમસ્યા ઊભી કરે છે. ``જો \dots તો \dots'' એવાં કેટલાંક વાક્યો
સંજોગાધીન રીતે સત્ય અથવા સંજોગાધીન રીતે અસત્ય હોય છે જ---હકીકતમાં,
સામાન્ય રીતે તે ન તો અનિવાર્ય હોય છે, ન તો અશક્ય.

\begin{prob}
  નીચેના \Log{S5}-પ્રમાણો આપો:
  \begin{enumerate}
    \item  $\Box \lnot !A \Entails !A \strictif !B$
    \item $!A \strictif !B \Entails \Box(!A \strictif !B)$
    \item $\lnot(!A \strictif !B)  \Entails \Box\lnot(!A \strictif !B)$
  \end{enumerate}
  તેના માટે $\strictif$ની વ્યાખ્યા વાપરો.
\end{prob}

\end{document}
